\documentclass[11pt,a4paper]{article}
\usepackage[margin=1in]{geometry}
\usepackage{amsmath,amssymb}
\usepackage{parskip}
\usepackage{hyperref}

\title{\textbf{Pre-Experiment Overview:\\ Determination of the $e/m$ Ratio of the Electron}}
\author{IBRAHIM H.I. ABUSHAWISH \\ Istanbul University, Department of Physics}
\date{}

\begin{document}
\maketitle

\section*{Purpose of the Experiment}
The main mission here is to pin down the $e/m$ ratio---the charge-to-mass ratio---of an electron. We pull this off by studying how cathode rays (which are just a beam of free electrons) behave after being accelerated by an electric potential and fired directly into a uniform magnetic field.

\section*{Importance and Significance}
This experiment is awesome because it lets us calculate the fundamental, invisible properties of an electron just by forcing it to steer in a circle. First, we heat a cathode to release electrons and accelerate them using a specific voltage, $V$, which turns their electrical energy into kinetic energy. Then, we power up a pair of Helmholtz coils to create a uniform magnetic field, $B$, in the space the electrons are flying through.

When those fast-moving electrons hit the magnetic field, they experience a magnetic force that pushes perpendicular to their velocity. This perpendicular push acts as a perfect centripetal force, bending the straight electron beam into a continuous circular orbit. By measuring the radius, $r$, of this glowing circular path on a fluorescent grid screen using $x$ and $y$ coordinates, we get our final variable. Plugging the accelerating voltage, the magnetic field strength, and the orbit's radius into the formula $\frac{e}{m} = \frac{2V}{B^2 r^2}$, we can calculate the exact charge-to-mass ratio of the electron. It's basically tracking the exact geometry of a curved trajectory to reverse-engineer the physics of the object itself!

\end{document}
